\documentclass[10pt,a4paper]{amsart}
\usepackage[margin=19mm]{geometry}
\usepackage{amsmath,amssymb,mathtools}
\usepackage[scr=rsfs,cal=euler]{mathalfa}
\usepackage{xfrac}
\usepackage[unicode,hidelinks,bookmarks=false]{hyperref}
\newcommand{\ie}{i.e.\ }
\newcommand{\cf}{cf.\ }
\newcommand{\resp}{respectively\ }
\newcommand{\Subsection}{\subsection}
\providecommand{\nobreakdash}{\nobreak\hbox{-}}
\setlength{\emergencystretch}{2em}
\multlinegap0pt
\usepackage{amssymb}
\usepackage{xcolor}
\usepackage{tikz}
\newtheorem{prop}{Proposition}
\newcommand{\I}{\mathcal{I}}
\renewcommand{\P}{\mathcal{P}}
\newcommand{\Sat}{\mathbf{Sat}}
\renewcommand{\max}{\mathbf{max}}
\newcommand{\mdd}{\mathbf{mdd}}
\title{On minimal pattern-containing inversion sequences}
\author{Benjamin Testart}
\date{Source: arXiv:2602.12130v2 (CC BY 4.0)}
\begin{document}
\maketitle

\noindent\emph{Source and licence.} On minimal pattern-containing inversion sequences. Version arXiv:2602.12130v2 (CC BY 4.0), distributed by arXiv under the Creative Commons Attribution 4.0 International licence, http://creativecommons.org/licenses/by/4.0/, as recorded in the arXiv licence metadata for this exact version and shown on the abstract page https://arxiv.org/abs/2602.12130v2. Full licence notice: \texttt{licenses/arxiv-2602.12130-CC-BY-4.0.txt}.\par\smallskip

\noindent\textit{Editorial scope.} Counted unit: the article's prop prop\_min\_max\_length (source0.tex lines 328--330). The article's complete original proof of it is delivered with it (source0.tex lines 331--375, one \texttt{proof} environment of 45 lines). No mathematics is written, repaired or re-derived: the statement and the proof are the article's own lines, in the article's own notation, and no retained line is substituted or re-flowed.  Retained uncounted as the source-local context the proof needs: lines 214--220; lines 313--315.  Nothing was omitted from any retained span, so the package carries no omission record.  No retained line contains a citation, so the wrapper carries no bibliography and no bibliography entry.  No retained line contains a figure environment, an image-inclusion command or drawing code, so no image is delivered and the package declares no asset dependency.  Editorial formatting only, in addition: an AMS \texttt{amsart} preamble that restates the article's own declarations and macros used by the retained lines, namely the environments \texttt{prop} on separate counters, together with the standard packages those lines need.  The retained environments print in this document as Proposition 1 in order of appearance, whereas the article prints the same items with the numbering of its own full document.  The complete native source of the article as distributed in the arXiv e-print for this exact version is bundled unchanged as \texttt{sources/194523/source0.tex}.
\begin{prop} \label{prop_minimal_characterization}
    Let $n,k \geq 1$, let $\rho \in \P_k$, and let $\sigma \in \I_n \cap \P_n$ be a sequence containing the pattern $\rho$. Then $\sigma$ is $\rho$-minimal if and only if for any set $R \subseteq [1,n]$ such that $\rho' := (\sigma_{i})_{i \in R}$ is an occurrence of $\rho$, both conditions below are satisfied:
    \begin{enumerate}
        \item $[\max(\Sat(\sigma)),n] \subseteq R$,
        \item any value of $\sigma$ which does not appear in $\rho'$ appears at least twice in $\sigma$.
    \end{enumerate}
\end{prop}

\begin{prop} \label{prop_tight_ub}
    Let $k \geq 2$, let $\rho \in \P_k$, and let $m := \mdd(\rho)$. If $\rho_i = 0$ for some $i \geq 2$, then there exists a $\rho$-minimal inversion sequence of length $k + 2m$.
\end{prop}

\begin{prop} \label{prop_min_max_length}
    For any pattern $\rho$, there exists a $\rho$-minimal inversion sequence of length at least $|\rho| + 2 \, \mdd(\rho) - 2$.
\end{prop}

\begin{proof}
    Let $k := |\rho|$, and let $m := \mdd(\rho)$. If $m \leq 2$, then $0^m \cdot \rho$ is a $\rho$-minimal inversion sequence of length $k + m \geq k + 2m - 2$. Assume now that $m \geq 3$. If $\rho_i = 0$ for some $i \geq 2$, then there exists a $\rho$-minimal inversion sequence of length $k + 2m$ by Proposition \ref{prop_tight_ub}. Assume now that $\rho_1$ is the only zero of $\rho$.

    Assume that $\rho_2 \geq 2$. Let $\rho' := (\rho_i + m - 1)_{i \in[2,k]}$, which is a sequence such that $\mdd(\rho') = 2m$, and let $\tau := 001122 \dots (m-1) (m-1) \cdot \rho' \in \I_{k+2m-1}$. Every term of value less than $m$ in $\tau$ can only take the role of the 0 in an occurrence of $\rho$, so all occurrences of $\rho$ in $\tau$ are subsequences of the form $v \cdot \rho'$ for some $v \in [0,m-1]$. These occurrences all satisfy the conditions of Proposition \ref{prop_minimal_characterization}, hence $\tau$ is $\rho$-minimal.

    Assume that $\rho_2 = 1$, and let $\ell := \max(i \in [1,k] \; : \; (\rho_j)_{j \in [1,i]} \in \I)$ be the length of the longest prefix of $\rho$ that is an inversion sequence. In particular, $\ell \geq 2$, $\rho_i < \ell$ for all $i \in [1,\ell]$, $\rho_{\ell+1} > \ell$, and there exists some $q \in [\ell+2,k]$ such that $\rho_q = \ell$.
    Let $\alpha := (\rho_i)_{i \in [1,\ell-1]}$, let $\beta := (\ell+1)(\ell+1)(\ell+2)(\ell+2) \dots (\ell+m-2)(\ell+m-2)$, and let $\gamma$ be the sequence of length $k-\ell+1$ defined by
    $$\forall i \in [\ell,k], \; \gamma_{i-\ell+1} := \begin{cases}
        \rho_{i} & \text{ if} \quad \rho_{i} \leq \ell, \\
        \rho_{i} + m-2 & \text{ if} \quad \rho_{i} > \ell.
    \end{cases}$$
    Let $\sigma := 00 \cdot \alpha \cdot \beta \cdot \gamma$, 
    which is a sequence of length $n := k+2m-2$. For example, if $\rho = 01165342$, then we have $m = 3$, $\ell = 3$, $\alpha = 01$, $\beta = 44$, and $\gamma = 176352$, hence $\sigma = 000144176352$.
    Going back to the general case, note that $\mdd(\alpha) = 0$, $\mdd(\beta) = \ell+1 = |00 \cdot\alpha|$, and $$\mdd(\gamma) = m+(\ell-1)+(m-2) = 2m+\ell-3 = |00 \cdot \alpha \cdot \beta|,$$ therefore $\sigma$ is an inversion sequence.
    It can be seen that $\sigma$ is also a Cayley permutation:
    \begin{itemize}
        \item every value in $[0, \ell]$ appears in $\rho$, and remains unchanged in $\alpha \cdot \gamma$,
        \item every value in $[\ell+1, \ell+m-2]$ appears in $\beta$,
        \item every value in $[\ell+m-1, \max(\rho)+m-2]$ appears in $\gamma$,
        \item by construction, $\max(\sigma) = \max(\gamma) = \max(\rho)+m-2$.
    \end{itemize}
    Now, we show that in an occurrence of $\rho$ in $\sigma$, the values of $\beta$ are too large to take the role of the term $\rho_i$ for any $i \leq \ell$. Let $r_1 < \dots < r_k \in [1,n]$ be such that $(\sigma_{r_i})_{i \in [1,k]}$ is an occurrence of $\rho$. Notice that $r_2 \geq 4$, since $\sigma_1 = \sigma_2 = \sigma_3 = 0$ and $\rho_1$ is the only zero of $\rho$. This implies that $r_i \geq i+2$ for all $i \in [2,k]$, in particular $r_{\ell} \geq \ell+2$. Let $b := |00 \cdot \alpha \cdot \beta| = 2m+\ell-3$, and let $q \in [\ell+2,k]$ be such that $\rho_q = \ell$.
    \begin{itemize}
        \item Assume for the sake of contradiction that $r_\ell \leq b$ and $r_{\ell+1} > b$. In particular, $\sigma_{r_{\ell}} > \ell$ since $r_\ell \in [\ell+2, b]$. Because $\rho_q > \rho_\ell$, this also implies that $\sigma_{r_q} > \ell$. Now, we have
        \begin{itemize}
            \item $\{i \; : \; \rho_i \geq \ell \} = \{i \; : \; \sigma_{r_i} \geq \sigma_{r_q}\}$ by definition of $(r_i)_{i \in [1,k]}$ and because $\rho_q = \ell$,
            \item $|\{i \; : \; \sigma_{r_i} \geq \sigma_{r_q}\}| \leq |\{i \; : \; \gamma_i \geq \sigma_{r_q}\}|$ since 
            \begin{itemize}
                \item $\sigma_{r_i} \geq \sigma_{r_q}$ is equivalent to $\rho_i \geq \ell$, so it implies that $i > \ell$ by definition of $\ell$,
                \item $r_{\ell+1} > b$ by hypothesis, so $(\sigma_{r_i})_{i \in [\ell+1,k]}$ is a subsequence of $\gamma$,
            \end{itemize}
            \item $|\{i \; : \; \gamma_i \geq \sigma_{r_q}\}| \leq |\{i \; : \; \gamma_i > \ell\}|$ since $\ell < \sigma_{r_q}$,
            \item $|\{i \; : \; \gamma_i > \ell\}| = |\{i \; : \; \rho_i > \ell\}|$ by definition of $\gamma$.
        \end{itemize}
        In summary, we have $|\{i \; : \; \rho_i \geq \ell \}| \leq |\{i \; : \; \rho_i > \ell\}|$; this implies that $\rho$ does not contain the value $\ell$, a contradiction.
        \item Assume for the sake of contradiction that $r_{\ell+1} \leq b$. Recall that $\sigma_{r_\ell} < \sigma_{r_q} < \sigma_{r_{\ell+1}}$, and $\ell + 2 \leq r_\ell < r_{\ell+1} < r_q$. Since $\ell+2 \leq r_\ell < r_{\ell+1} \leq b$ and $\beta$ is nondecreasing, we have $\ell+1 \leq \sigma_{r_\ell} \leq \sigma_{r_{\ell+1}} \leq \ell + m - 2$. By construction of $\sigma$, every term $\sigma_i$ for $i \in [1+r_{\ell+1}, n]$ falls in one of three cases:
        \begin{itemize}
            \item $\sigma_i \geq \sigma_{r_{\ell+1}}$ if $i \leq b$ since $\beta$ is nondecreasing,
            \item $\sigma_i \leq \ell < \sigma_{r_\ell}$ if $i > b$ and $\rho_{i+k-n} \leq \ell$,
            \item $\sigma_i > \ell+m-2 \geq \sigma_{r_{\ell+1}}$ if $i > b$ and $\rho_{i+k-n} > \ell$.
        \end{itemize}
        This contradicts the existence of $\sigma_{r_q}$.
    \end{itemize}
    We showed that $r_\ell > b$. Since $|[b+1,n]| = |[\ell,k]|$, we have ${r_i = i+b-\ell+1 = i+n-k}$ for all $i \in [\ell,k]$. In particular, $\sigma_{r_q} = \sigma_{q+b-\ell+1} = \gamma_{q-\ell+1} = \rho_q = \ell$, hence ${\sigma_{r_{\ell-1}} < \ell}$ since ${\rho_{\ell-1} < \rho_q}$. This implies that $r_{\ell-1} \leq \ell+1$, since $r_{\ell-1} < r_\ell = b+1$ and $\sigma_i > \ell$ for all ${i \in [\ell+2,b]}$. Lastly, recalling that $r_i \geq i+2$ for all $i \geq 2$, we have ${r_j = j+2}$ for all $j \in [2,\ell-1]$. In summary, there are at most three occurrences of $\rho$ in $\sigma$, as we completely determined the value of each $r_i$ except for $r_1 \in \{1,2,3\}$. It is easy to verify that those three subsequences are occurrences of $\rho$: indeed, all three are equal to the sequence $\alpha \cdot \gamma$, which is order-isomorphic to $\rho$ by definition. It is also straightforward to verify that all three occurrences of $\rho$ satisfy the conditions of Proposition \ref{prop_minimal_characterization}, since ${\max(\Sat(\sigma)) = b+\max(\Sat(\rho)) \geq b+2}$, and the entries of $\sigma$ which are not used in each occurrence of $\rho$ are exactly two zeros and the factor $\beta$. We conclude that $\sigma$ is $\rho$-minimal.
\end{proof}

\end{document}
